%% =====================================================================
%%  Bilan 4 — Maintenant et bientôt (unités 13 à 16)
%%  Pistes : B4.01 (écoute notée) · B4.02 (expression orale)
%% =====================================================================
\bilan{4}

\noindent{\fontsize{11}{14.6}\selectfont Quatrième partie terminée ! Vous savez maintenant dire ce
qui se passe, parler du temps qu'il fait et des dates, inviter quelqu'un et présenter vos projets.
Révisez l'essentiel, puis faites les huit exercices notés sur 20. Les corrigés sont
page~\pageref{cor-b4-1}.\par}

%% ---------------------------------------------------------------------
\rubrique{grammaire}{Révisez}

\begin{paire}
\begin{regle}[Le présent en \en{-ing} · unité 13]
\en{I'm working.} · \en{She isn't working.} · \en{Are you working?}\par\vspace{1mm}
\en{make} → \en{making} · \en{sit} → \en{sitting}\par\vspace{1mm}
{\fontsize{9.6}{12}\selectfont Habitude : \en{He usually wears jeans.} Maintenant : \en{Today, he's
wearing a suit.}\par}
\end{regle}
\begin{regle}[Les dates · unité 14]
\en{first}, \en{second}, \en{third}, \en{fourth}… \en{twentieth}, \en{twenty-first}\par\vspace{1mm}
\en{at seven} · \en{on Monday}, \en{on March 15th} · \en{in May}, \en{in winter}\par\vspace{1mm}
\en{What's the weather like? — It's raining.}
\end{regle}
\end{paire}

\begin{paire}
\begin{regle}[Inviter · unité 15]
\en{Let's go out!} · \en{Would you like to come?}\par
\en{I'd love to!} · \en{I'm sorry, I can't.}\par\vspace{1mm}
Pronoms compléments : \en{me}, \en{you}, \en{him}, \en{her}, \en{it}, \en{us}, \en{them}
\end{regle}
\begin{regle}[\en{be going to} · unité 16]
\en{I'm going to travel.} · \en{He isn't going to fly.}\par
\en{Are you going to stay? — Yes, I am.}\par\vspace{1mm}
\en{tomorrow} · \en{next week} · \en{in two weeks} · \en{soon}
\end{regle}
\end{paire}

%% ---------------------------------------------------------------------
\rubrique{pratique}{Testez-vous}

\exercice[ecrire,niveau=1,points=2]{La forme en \en{-ing}}{Écrivez la forme en \en{-ing}.
\textit{(0,5 point par réponse)}}
\begin{anglais}
\begin{questions}[4]
\item sit → \trou[20mm]{sitting}
\item write → \trou[20mm]{writing}
\item swim → \trou[20mm]{swimming}
\item play → \trou[20mm]{playing}
\end{questions}
\end{anglais}

\exercice[ecrire,niveau=2,points=3]{Maintenant ou d'habitude ?}{Mettez le verbe au présent simple ou
au présent en \en{-ing}. \textit{(0,5 point par réponse)}}
\begin{questions}
\item Look! Jack \trou[32mm]{is dancing}! \hfill (\en{dance})
\item Emma usually \trou[32mm]{drinks} tea. \hfill (\en{drink})
\item Shh! We \trou[32mm]{are watching} a movie. \hfill (\en{watch})
\item Sam \trou[32mm]{plays} the guitar every evening. \hfill (\en{play})
\item What \trou[32mm]{are you doing}? \hfill (\en{you / do})
\item \trou[32mm]{Does she work} on Saturdays? \hfill (\en{she / work})
\end{questions}

\exercice[ecrire,niveau=1,points=2]{\en{in}, \en{on} ou \en{at} ?}{Complétez. \textit{(0,5 point par
réponse)}}
\begin{anglais}
\begin{questions}[2]
\item See you \trou[12mm]{on} Monday!
\item It's hot \trou[12mm]{in} July.
\item The movie starts \trou[12mm]{at} six o'clock.
\item Christmas is \trou[12mm]{on} December 25th.
\end{questions}
\end{anglais}

\exercice[ecrire,niveau=2,points=3]{Les pronoms compléments}{Complétez avec le pronom complément qui
correspond aux mots entre parenthèses. \textit{(0,5 point par réponse)}}
\begin{anglais}
\begin{questions}[2]
\item Call \trou[14mm]{me} tomorrow! \hfill (I)
\item I don't know \trou[14mm]{him}. \hfill (Jack)
\item Come with \trou[14mm]{us}! \hfill (we)
\item Look at \trou[14mm]{them}! \hfill (the kids)
\item I like \trou[14mm]{it}. \hfill (the hat)
\item I love \trou[14mm]{her}. \hfill (Léa)
\end{questions}
\end{anglais}

\exercice[ecrire,niveau=2,points=2]{Projets}{Complétez avec \en{be going to} et le verbe.
\textit{(0,5 point par réponse)}}
\begin{questions}
\item Next week, my sister \trou[44mm]{is going to visit} us. \hfill (\en{visit})
\item Look at the clouds! It \trou[44mm]{is going to rain}. \hfill (\en{rain})
\item They \trou[44mm]{aren't going to fly}. \hfill (\en{not fly})
\item \trou[44mm]{Is she going to come} tonight? \hfill (\en{she / come})
\end{questions}

\exercice[ecouter,niveau=2,points=3,piste=B4.01]{Au téléphone}{Jack appelle Léa. Écoutez, puis indiquez
si chaque phrase est vraie (V) ou fausse (F). \textit{(0,5 point par réponse)}}
\begin{anglais}
\begin{questions}[2]
\item Léa is cooking dinner. \hfill\trou[9mm]{V}
\item Jack is playing basketball. \hfill\trou[9mm]{F}
\item Jack and Kate are going to go ice skating. \hfill\trou[9mm]{V}
\item They're going to meet on Sunday. \hfill\trou[9mm]{F}
\item They're going to meet at two. \hfill\trou[9mm]{V}
\item It's going to be warm. \hfill\trou[9mm]{F}
\end{questions}
\end{anglais}

\exercice[lire,niveau=2,points=3]{La carte postale d'Emma}{Lisez la carte postale d'Emma, puis indiquez
si chaque phrase est vraie (V) ou fausse (F). \textit{(0,5 point par réponse)}}
\begin{anglais}
\begin{tcolorbox}[enhanced,width=140mm,center,colback=white,colframe=otRule,boxrule=0.6pt,arc=1mm,
  left=5mm,right=5mm,top=4mm,bottom=4mm,drop shadow={black!25,opacity=0.5}]
\begin{minipage}[t]{0.6\linewidth}\vspace{0pt}
\fontsize{10.4}{14}\selectfont\itshape
Hi everyone!\par
I'm writing from London. It's raining (of course!) and it's cold: four degrees Celsius. Right now,
I'm sitting in a café with my sister. Tomorrow, we're going to visit the British Museum, and on
Saturday we're going to see a musical. I'm going to fly back on January 4th.\par
See you soon!\quad Emma
\end{minipage}\hfill{\color{otRule}\vrule width 0.5pt}\hfill
\begin{minipage}[t]{0.32\linewidth}\vspace{0pt}
\raggedleft
\begin{tikzpicture}
  \draw[otRedDark,line width=0.8pt,dash pattern=on 1.4pt off 1pt] (0,0) rectangle (16mm,19mm);
  \fill[otRedTint] (1mm,1mm) rectangle (15mm,18mm);
  \otDessinIcone{etoile}{otRedDark}{otRedTint}{8mm}{(8mm,11mm)}
  \node[font=\otSemi\fontsize{5.6}{7}\selectfont,text=otRedDark] at (8mm,4mm) {UK};
\end{tikzpicture}\par\vspace{5mm}
\raggedright\fontsize{9.6}{14}\selectfont
Léa Martin\par\rule{\linewidth}{0.4pt}\par
Bright Ideas\par\rule{\linewidth}{0.4pt}\par
New York, NY\par\rule{\linewidth}{0.4pt}\par
USA
\end{minipage}
\end{tcolorbox}
\begin{questions}[2]
\item Emma is in London. \hfill\trou[9mm]{V}
\item It's sunny. \hfill\trou[9mm]{F}
\item Emma is sitting in a café. \hfill\trou[9mm]{V}
\item Tomorrow, she's going to see a musical. \hfill\trou[9mm]{F}
\item She's with her sister. \hfill\trou[9mm]{V}
\item She's going to fly back in February. \hfill\trou[9mm]{F}
\end{questions}
\end{anglais}

\exercice[ecrire,niveau=3,points=2]{Votre week-end}{En quatre phrases, dites ce que vous allez faire ce
week-end, avec qui et quel temps il va faire. \textit{(0,5 point par phrase correcte)}}
\lignes{4}

%% ---------------------------------------------------------------------
\rubrique{ecouter}{Parlez}

\exercice[ecouter,parler,niveau=3,piste=B4.02]{Un appel d'Emma}{Emma vous appelle de Londres.
Répondez-lui \textbf{à voix haute} pendant chaque pause. Cet exercice n'est pas noté : comparez vos
réponses avec les modèles du corrigé.}
\begin{center}\otTableau
\begin{tabular}{@{}>{\otSemi\fontsize{8.6}{10}\selectfont\color{otNavy}}l@{\hspace{3mm}}>{\enOn}p{70mm}<{\enOff}@{\hspace{6mm}}p{78mm}@{}}
 & \otTH{Emma dit\ldots} & \otTH{Vous\ldots}\\\midrule
1. & Hi! What are you doing right now? & dites ce que vous faites.\\
2. & What's the weather like where you are? & décrivez le temps qu'il fait.\\
3. & When's your birthday? & donnez la date.\\
4. & What are you going to do this weekend? & présentez un projet.\\
5. & Would you like to go to the movies on Friday? & acceptez ou refusez poliment.\\\bottomrule
\end{tabular}\end{center}

\otScore{20}
